\documentclass[../DoAn.tex]{subfiles}
\begin{document}

\section{Tổng quan các đóng góp}
\label{section:5.0}

Chương này phân tích chi tiết sáu đóng góp kỹ thuật và khoa học của đồ án, cách chúng giải quyết các vấn đề cụ thể trong ứng dụng \gls{GraphRAG} cho miền pháp luật lao động Việt Nam, và giá trị tái sử dụng trong các ứng dụng pháp lý khác. Mỗi đóng góp được trình bày theo cấu trúc: \textbf{giá trị lõi} → \textbf{vấn đề được giải quyết} → \textbf{so sánh với công trình liên quan} → \textbf{minh chứng thực nghiệm} (tham chiếu Chương 4). Các đóng góp bổ sung lẫn nhau, từng phần giải quyết một khía cạnh khác của pipeline Legal GraphRAG.

\section{Ontology hai lớp với canonicalization để bảo toàn cấu trúc pháp lý}
\label{section:5.1}

\subsection{Giá trị lõi và vấn đề đặt ra}
\label{subsection:5.1.1}

Vấn đề cốt lõi khi áp dụng \gls{RAG} cho pháp luật là \textbf{mất cấu trúc pháp lý} khi văn bản bị chunking tùy ý:

\begin{itemize}[nosep, leftmargin=*]
    \item \textbf{Naive RAG:} Chunk theo token cố định (512--1024 token, overlap 50\%) → cắt giữa Khoản, mất bối cảnh pháp lý. Trích dẫn không thể trace về đơn vị chuẩn (Điều/Khoản), gây mất tin cậy trong domain pháp luật.
    
    \item \textbf{GraphRAG gốc:} Trích entity/quan hệ bằng LLM nhưng schema mặc định (ORGANIZATION, PERSON, LOCATION) không phù hợp domain luật lao động Việt Nam. Thiếu lớp cấu trúc văn bản deterministic (Văn bản pháp lý $\rightarrow$ Phần $\rightarrow$ Chương $\rightarrow$ Điều $\rightarrow$ Khoản) → không đảm bảo nhất quán đồ thị.
    
    \item \textbf{Hậu quả:} Không suy luận multi-hop giữa các quy phạm, không trích dẫn chính xác, khó giải thích nguồn.
\end{itemize}

\subsection{Giải pháp: Ontology hai lớp + canonicalization}
\label{subsection:5.1.2}

Đóng góp cốt lõi là thiết kế \textbf{ontology hai lớp với canonicalization} (chi tiết Mục~\ref{subsection:3.2.2}):

\textbf{Lớp 1 --- Cấu trúc văn bản (100\% deterministic):} Biểu diễn phân cấp hình thức (Văn bản pháp lý, Phần, Chương, Mục, Điều, Khoản, Điểm) và các quan hệ phân cấp (\texttt{contains}, \texttt{nested\_in}), tham chiếu chéo (\texttt{cites}, \texttt{amends}, \texttt{repeals}, \texttt{guided\_by}). Mỗi thực thể gán định danh chuẩn cố định (ví dụ: \texttt{BLLĐ\_45\_2019\_Điều\_90\_Khoản\_1}), không dùng LLM, độ tin cậy 100\%, lặp lại được.

\textbf{Lớp 2 --- Ngữ nghĩa pháp luật (LLM extraction):} Trích xuất 14 loại thực thể domain (ChuThe, HanhVi, CoQuan, TienLuong, NghiPhep, CheTai, v.v.) và các quan hệ logic pháp lý (\texttt{entitles}, \texttt{obligates}, \texttt{prohibits}, \texttt{penalizes}). Schema entity types được khai báo tường minh, không sử dụng schema mặc định của GraphRAG gốc.

\textbf{Lớp 1.5 --- Canonicalization:} Ánh xạ các biểu diễn khác nhau (``Điều 35'', ``khoản 1 Điều 35 BLLĐ 2019'', ``mục 1 Điều 35'') về cùng một định danh chuẩn. Giải quyết xung đột giữa output LLM (chuỗi tự do) và ID cấu trúc (ID cố định) khi hợp nhất, tránh trùng lặp node (cùng Điều 35 nhưng hai ID khác nhau).

Lựa chọn thiết kế: \textbf{không extract Lớp 1 bằng LLM} --- vì LLM không đảm bảo nhất quán giữa các lần chạy, dẫn tới duplicate node hay ngộ nhận giữa Điều khác nhau. Điều này là điểm khác biệt quan trọng so với GraphRAG gốc, nơi toàn bộ extraction dùng LLM.

\subsection{So sánh với công trình liên quan}
\label{subsection:5.1.3}

\begin{itemize}[nosep, leftmargin=*]
    \item \textbf{Legal KG tay-craft (LKIF, Akoma Ntoso):} Xây dựng thủ công từng ontology, độ chính xác cao nhưng tốn công bảo trì, khó mở rộng khi luật thay đổi.
    
    \item \textbf{Giải pháp của đồ án:} Kết hợp parser deterministic (L1) + LLM (L2), cho phép tự động cập nhật khi thêm văn bản mới. L1 đảm bảo nhất quán; L2 bổ sung ngữ cảnh ngữ nghĩa.
    
    \item \textbf{GraphRAG gốc:} Toàn bộ LLM → linh hoạt nhưng có rủi ro không nhất quán. Ontology của đồ án thêm lớp cấu trúc xác định, cải thiện tin cậy.
\end{itemize}

\subsection{Minh chứng thực nghiệm}
\label{subsection:5.1.4}

\begin{itemize}[nosep, leftmargin=*]
    \item Parse thành công 581 Điều, 2069 Khoản từ 7 văn bản (độ chính xác 100\% trên cấu trúc hình thức).
    
    \item Đồ thị merged: $\approx$ 5120 entity (L1 $\approx$ 2800 + L2 $\approx$ 2320), $\approx$ 7890 quan hệ, 15 community.
    
    \item Chunking theo Điều (overlap=0) → Local Search đạt \textbf{100\%} keyword accuracy trên nhóm \texttt{single\_hop} (Bảng~\ref{table:by_category}, Chương~\ref{chapter:Experiment}).
    
    \item Citation accuracy (trích dẫn đúng Điều/Khoản): \textbf{88\%} trên Local Search --- cao hơn Naive RAG ($\approx$ 60--70\% vì cắt ngang Khoản).
    
    \item \textbf{Reusable framework:} Pattern L1 (parser) + L2 (LLM extraction) có thể mở rộng sang luật dân sự, hình sự bằng cách thay parser regex và khai báo entity types riêng --- không cần thiết kế lại toàn bộ ontology.
\end{itemize}

\section{Multi-hop Reasoning Engine để suy luận chéo quy phạm}
\label{section:5.2}

\subsection{Vấn đề và giá trị}
\label{subsection:5.2.1}

Nhiều câu hỏi pháp lý không thể trả lời bằng một Điều duy nhất. Chúng yêu cầu suy luận qua chuỗi quan hệ:

\begin{example}
``NSDLĐ không trả trợ cấp thôi việc sẽ bị xử lý thế nào?''
Trace: Điều 46 BLLĐ (nghĩa vụ trả trợ cấp) $\rightarrow$ NĐ 12/2022 Điều 3 (chế tài vi phạm) $\rightarrow$ Mức phạt tiền.
\end{example}

Vấn đề: \textbf{Local Search} retrieve text units ngữ nghĩa gần nhau nhưng không đảm bảo retrieve \emph{cả hai} văn bản (BLLĐ \textbf{và} Nghị định) trong chuỗi. Thí nghiệm: Local đạt 33\% trên nhóm \texttt{multi\_hop} (LD010--LD012), chứng minh vector similarity đơn thuần không đủ.

Giá trị: \textbf{Động cơ suy luận logic pháp lý}, không dựa trên ngữ nghĩa mà dựa trên \textbf{quan hệ pháp lý} (``hành vi A vi phạm → bị phạt theo Điều B → mức phạt là X'').

\subsection{Giải pháp: VNLegalReasoningEngine}
\label{subsection:5.2.2}

Triển khai engine trên merged graph với ba chain templates:

\begin{itemize}[nosep, leftmargin=*]
    \item \textbf{Violation chain:} Hành vi vi phạm $\rightarrow$ Chế tài. Quan hệ: \texttt{HanhVi} --\texttt{penalizes}$\rightarrow$ \texttt{CheTai} --\texttt{cites}$\rightarrow$ \texttt{Dieu}.
    
    \item \textbf{Entitlement chain:} Chủ thể $\rightarrow$ Quyền lợi $\rightarrow$ Điều kiện. Quan hệ: \texttt{ChuThe} --\texttt{entitles}$\rightarrow$ \texttt{QuyenLoi} --\texttt{applies\_to}$\rightarrow$ \texttt{DieuKien}.
    
    \item \textbf{Procedure chain:} Hành động $\rightarrow$ Quy trình $\rightarrow$ Cơ quan. Quan hệ: --\texttt{applies\_to}$\rightarrow$ --\texttt{enforced\_by}$\rightarrow$.
\end{itemize}

Thuật toán: DFS trên graph; độ sâu tối đa 4 hop để tránh explosion. Output: \texttt{ReasoningChain} gồm chuỗi \texttt{steps[]}, danh sách \texttt{cited\_articles[]}, câu trả lời \texttt{final\_answer}.

\subsection{So sánh với công trình liên quan}
\label{subsection:5.2.3}

\begin{itemize}[nosep, leftmargin=*]
    \item \textbf{Temporal GraphRAG:} Multi-hop trên timeline sự kiện tin tức → không phù hợp pháp luật (không có sự kiện, chỉ có quan hệ logic).
    
    \item \textbf{Retrieval-based multi-hop (HyDE, etc.):} Retrieve nhiều lần từng bước → chậm, cộng dồn lỗi mỗi bước.
    
    \item \textbf{Giải pháp của đồ án:} Duyệt trực tiếp trên graph --- nhanh ($\approx$ 200--500ms), logic rõ ràng (trace quan hệ), lỗi từ extraction LLM được hạn chế (chỉ ảnh hưởng L2, L1 không sai).
\end{itemize}

\subsection{Minh chứng thực nghiệm}
\label{subsection:5.2.4}

\begin{itemize}[nosep, leftmargin=*]
    \item Local Search: 33\% keyword accuracy trên multi-hop. Local + Multi-hop: \textbf{100\%}.
    
    \item Case LD012: Trace violation từ ``lương dưới tối thiểu'' $\rightarrow$ NĐ 12/2022 --- chain\_hit đúng.
    
    \item Latency: $\approx$ 3.5s P50 tổng (Local ~3s + engine ~200--500ms). Chấp nhận được cho chat.
    
    \item Độ dài chain: 3--4 hop trung bình; thường kết thúc tại Điều cấu trúc (L1) hoặc chế tài (L2).
\end{itemize}

\section{Rule Validator VR001--VR016 cho phân tích HĐLĐ deterministic + hybrid}
\label{section:5.3}

\subsection{Vấn đề: Phân tích HĐLĐ LLM-only không tin cậy}
\label{subsection:5.3.1}

Phân tích hợp đồng lao động yêu cầu: (1) \textbf{Nhanh} (50--150 khoản); (2) \textbf{Chính xác} (không miss vi phạm); (3) \textbf{Giải thích được} (trích dẫn căn cứ).

LLM-only có hạn chế: chậm (1--2s/khoản), hallucination (trích sai Điều), chi phí token cao.

Quan sát khóa: Domain luật lao động VN có \textbf{nhiều ngưỡng số học rõ ràng} --- 12 ngày phép/năm, 60 ngày thử việc, 48 giờ/tuần, lương tối thiểu vùng I--IV. Các ngưỡng này phù hợp kiểm tra bằng \textbf{rule deterministic}, không cần LLM.

\subsection{Giải pháp: Kiến trúc hybrid ba lớp}
\label{subsection:5.3.2}

Thiết kế layer stack:

\textbf{Lớp 1 --- Rule-based} ($<$ 5ms/khoản): 16 quy tắc VR001--VR016 với ngưỡng cụ thể, regex pattern, kèm căn cứ pháp lý. Ví dụ:
\begin{itemize}[nosep, leftmargin=*]
    \item VR001: Lương monthly $<$ \texttt{REGIONAL\_MINIMUM\_WAGE[region]} → VIOLATION (Điều 90, NĐ 74/2024).
    \item VR008: Nghỉ phép năm $<$ 12 ngày → VIOLATION (Điều 113).
\end{itemize}

\textbf{Lớp 2 --- Graph grounding} (200--500ms/khoản phức tạp): Chỉ khoản category phức tạp hoặc UNKNOWN mới map qua \texttt{VNLegalReasoningEngine} → \texttt{MappedLawSnippet}. Cung cấp context pháp lý cho LLM mà không thay thế rule.

\textbf{Lớp 3 --- LLM bổ sung} (batch processing): Phân tích rủi ro ngữ nghĩa (ví dụ: lương thêm, phụ cấp, điều kiện riêng) trên khoản chưa có VIOLATION rõ từ VR001--VR007. Skip LLM nếu rule đã báo rõ → tiết kiệm 80--90\% token.

\subsection{So sánh với công trình liên quan}
\label{subsection:5.3.3}

\begin{itemize}[nosep, leftmargin=*]
    \item \textbf{LLM-only contract analysis:} GPT-4 detect violations --- chậm, bỏ qua chi tiết, không giải thích đủ.
    
    \item \textbf{Rule-only:} Kiểm tra deterministic --- nhanh nhưng không nắm rủi ro ngữ nghĩa (``cách trả lương bất lợi NLĐ'').
    
    \item \textbf{Giải pháp của đồ án:} Hybrid --- rule xử lý 95\% khoản ($<$ 5ms), LLM xử lý 5\% khoản phức tạp. Kết hợp graph grounding (L2) để cung cấp context mà không phụ thuộc 100\% vào LLM.
\end{itemize}

\subsection{Minh chứng thực nghiệm}
\label{subsection:5.3.4}

Case study HĐLĐ mẫu (7 khoản):
\begin{itemize}[nosep, leftmargin=*]
    \item VR008: Phát hiện đúng --- nghỉ phép 6 ngày/năm $<$ 12 → VIOLATION (Điều 113).
    \item VR015: Phát hiện --- thiếu địa điểm → MEDIUM\_RISK (Điều 21).
    \item L001 (LLM): Làm rõ rủi ro lương/phụ cấp → HIGH\_RISK.
    \item Điểm tuân thủ: 74/100; thời gian xử lý $\approx$ 25--40s (DOCX, 7 khoản) --- chấp nhận được.
    \item Trace giải thích: Mỗi issue ghi \texttt{rule\_id} (VR008) + căn cứ (Điều 113) → tin cậy.
\end{itemize}

VR validator là đóng góp \textbf{domain-specific}, chưa có trong GraphRAG gốc, kết hợp hybrid mà không thay thế rule.

\section{Temporal Validity Filter lọc trích dẫn lỗi thời}
\label{section:5.4}

\subsection{Vấn đề: Luật thay đổi liên tục}
\label{subsection:5.4.1}

Pháp luật lao động VN thay đổi liên tục (lương vùng, chế tài, v.v.). Trích dẫn \textbf{văn bản hết hiệu lực} là sai pháp lý nghiêm trọng (ví dụ: NĐ 90/2019 → thay bằng NĐ 74/2024).

GraphRAG gốc và RAG truyền thống không có cơ chế lọc theo \textbf{hiệu lực pháp lý tạm thời} (temporal validity của văn bản, không phải timeline sự kiện).

\subsection{Giải pháp: Temporal filter từ metadata}
\label{subsection:5.4.2}

Module \texttt{query/temporal\_filter.py}:

\begin{itemize}[nosep, leftmargin=*]
    \item Load \texttt{metadata.json}: mỗi văn bản có \texttt{tinh\_trang} (\texttt{con\_hieu\_luc}/\texttt{het\_hieu\_luc}), \texttt{ngay\_hieu\_luc}, \texttt{huong\_dan\_boi} (văn bản hướng dẫn).
    
    \item Sau retrieve, lọc \texttt{article\_citations} theo \texttt{is\_effective(van\_ban\_id, as\_of=today)}.
    
    \item Sinh \texttt{temporal\_warnings} nếu trích dẫn hết hiệu lực → thông báo người dùng.
\end{itemize}

Khác biệt vs. Temporal GraphRAG: Filter này xử lý \textbf{hiệu lực pháp lý văn bản} (metadata tĩnh), không phải timeline sự kiện bên ngoài.

\subsection{So sánh}
\label{subsection:5.4.3}

\begin{itemize}[nosep, leftmargin=*]
    \item \textbf{Temporal GraphRAG:} Suy luận theo thời gian sự kiện tin tức (``sự kiện X xảy ra trước Y'').
    
    \item \textbf{Giải pháp của đồ án:} Cảnh báo hiệu lực văn bản pháp lý --- tập trung vào bối cảnh pháp lý VN, nơi văn bản thường được thay thế (lương vùng hằng năm, chế tài cập nhật).
\end{itemize}

\subsection{Minh chứng}
\label{subsection:5.4.4}

Thí nghiệm E5 (TMP001--TMP002):
\begin{itemize}[nosep, leftmargin=*]
    \item Với filter: Trả NĐ 74/2024 (hiện tại); cảnh báo NĐ 90/2019 hết hiệu lực.
    
    \item Không filter: Vẫn có thể nhắc NĐ 90/2019 nếu còn text units.
    
    \item Nhóm \texttt{temporal}: Local + filter → \textbf{100\%} keyword accuracy.
\end{itemize}

\section{Kiến trúc hybrid hỏi đáp pháp luật: tích hợp end-to-end}
\label{section:5.5}

\subsection{Giá trị: Framework tái sử dụng cho Legal RAG}
\label{subsection:5.5.1}

Đóng góp thứ sáu là \textbf{kiến trúc end-to-end} tích hợp các thành phần trên thành hai luồng ứng dụng thực tiễn:

\textbf{Luồng Chat hỏi đáp pháp luật:}
\[
\text{Query} \xrightarrow{\text{Mode}} \text{Retrieve} \xrightarrow{\text{[+Multi-hop]}} \text{[+Temporal filter]} \xrightarrow{\text{LLM}} \text{Response + Citations + Warnings}
\]

\textbf{Luồng Phân tích HĐLĐ:}
\[
\text{Upload} \xrightarrow{\text{Segment}} \text{VR rules} \xrightarrow{\text{[+Graph]}} \text{[+LLM]} \xrightarrow{\text{Report}} \text{Tuân thủ \%}
\]

Nguyên tắc \textbf{triple-check}: (1) Retrieval grounding (corpus đã index); (2) Temporal check (hiệu lực); (3) Rule check (HĐLĐ, deterministic trước LLM).

Mọi output có \textbf{trace giải thích}: \texttt{article\_citations}, \texttt{reasoning\_chain}, \texttt{rule\_id} → tin cậy.

\subsection{So sánh với GraphRAG gốc}
\label{subsection:5.5.2}

\begin{itemize}[nosep, leftmargin=*]
    \item \textbf{GraphRAG gốc:} Local Search, Global Search (community map-reduce) --- tốt cho tin tức.
    
    \item \textbf{Đồ án:} Local + Global + Multi-hop + Temporal filter + Rule validator + Contract analyzer --- tối ưu hóa cho pháp luật.
    
    \item \textbf{Routing logic:} Chọn mode theo loại câu hỏi (tra cứu → Local; tổng quan → Global; suy luận → Multi-hop; lỗi thời → Temporal).
\end{itemize}

\subsection{Minh chứng tổng hợp}
\label{subsection:5.5.3}

\begin{itemize}[nosep, leftmargin=*]
    \item Domain \texttt{lao\_dong} (n=10): Local 90\%, Local + Multi-hop 100\% keyword accuracy.
    
    \item Single-hop: 100\% với ontology L1 + chunking theo Điều.
    
    \item Multi-hop: 33% → 100% nhờ engine.
    
    \item HĐLĐ case: VR008 phát hiện, Điều 113 → tin cậy.
    
    \item Pipeline end-to-end: indexing → query → API → UI --- demo hoạt động.
\end{itemize}

\section{Tổng kết các đóng góp}
\label{section:5.6}

Chương này đã phân tích sáu đóng góp chính:

\begin{enumerate}[nosep, leftmargin=*]
    \item \textbf{Ontology hai lớp + canonicalization} --- giải quyết mất cấu trúc pháp lý, nền tảng citation accuracy cao, reusable.
    
    \item \textbf{Multi-hop reasoning engine} --- suy luận chéo quy phạm, cải thiện 67 điểm phần trăm trên nhóm multi-hop.
    
    \item \textbf{Rule validator VR001--VR016} --- kiểm tra HĐLĐ nhanh, deterministic, hybrid (rule + graph + LLM).
    
    \item \textbf{Temporal validity filter} --- lọc văn bản hết hiệu lực, cảnh báo người dùng, phù hợp bối cảnh pháp lý VN.
    
    \item \textbf{Kiến trúc hybrid tích hợp} --- end-to-end routing, triple-check tin cậy, FastAPI + Next.js.
    
    \item \textbf{Framework tái sử dụng} --- pattern L1 + L2 + mở rộng sang luật dân sự, hình sự.
\end{enumerate}

Các đóng góp giải quyết các vấn đề cụ thể trong lĩnh vực pháp luật (trích dẫn chính xác, suy luận logic, hiệu lực tạm thời) mà GraphRAG gốc và RAG truyền thống không xử lý được. Framework được thiết kế để tái sử dụng trong các domain pháp lý khác bằng cách thay đổi parser, entity types, và rule validator mà không cần thiết kế lại toàn bộ kiến trúc. Chương~\ref{chapter:conclusion} tổng kết hạn chế và hướng phát triển; Chương~\ref{chapter:recommendations} đưa ra khuyến nghị triển khai thực tế.

\end{document}
